% Nouveau EGA V (cote 25): archive cote 25, site page 119 (PDF page 120).
% typescript. Draft compared: none.
% AI-transcribed by Claude Opus 5.5 (claude-opus-5-5, adaptive thinking, effort high), Message Batches API, batch ega25b (msgbatch_01A7n5YuXubUPNK1zsepi8RG), 2026-11-13; not yet checked by eye.
\scanpage{25/119}
\var{\margin{\unsure{2/2}}}{\placedmargin{\unsure{2/2}}}{marginal note placed}{ega25_25_119-v0001}
\var{\margin{(trait vertical au crayon dans la marge gauche, le long du premier paragraphe)}}{\placedmargin{(trait vertical au crayon dans la marge gauche, le long du premier paragraphe)}}{marginal note placed}{ega25_25_119-v0002}
Je me répens d'avoir donné la proposition \var{\deleted{\unsure{3}}}{}{deletion}{ega25_25_119-v0003} sous une forme éculée et
pisseuse, à mi chemin entre  les hypothèses classiques et les hypothèses
naturelles, et sans donner une récip\unsure{r}oque. Aussi je propose l'énoncé:
\uline{Proposition\unsure{\hand{3}}3. \var{\deleted{revue et corrigée}}{}{deletion}{ega25_25_119-v0004}}.  Soit $X \to S$ plat loc de présentation
finie,  $Q=\check{P} = P(\check{E})$  un fibré \var{\xout{\unsure{vectoriel sur}}}{}{typewriter strike-out}{ega25_25_119-v0005} projectif sur $S$ défini par
un Module loc  libre de type fini $F = \check{E}$ ,  considérons l'ensemble
$\emptyset$  des systêmes linéaires \var{\inserted{$D$ \var{\deleted{\unsure{E}}}{}{deletion}{ega25_25_119-v0007}}}{$D$ \var{\deleted{\unsure{E}}}{}{deletion}{ega25_25_119-v0007}}{insertion}{ega25_25_119-v0006} de diviseurs \var{\xout{\unsure{sur}}}{}{typewriter strike-out}{ega25_25_119-v0008} sur $X$, paramétrés par $Q$, tels
que l'ensemble  $Z$ des points fixes de $D$  satisfasse la propriété: $z \in Z$
implique prof$_z \underline{O}_{X_s} \geqslant 2$  (où  $s$ \var{\xout{\unsure{$= g(z)$)}}}{}{typewriter strike-out}{ega25_25_119-v0009} est l'image de $z$ dans $S$). \var{\xout{\unsure{Alors}}}{}{typewriter strike-out}{ega25_25_119-v0010}
\var{\xout{\unsure{$D \in \emptyset$ associons le morphisme correspondant $f$ $X$-$Z$ $P$ considérée}}}{}{typewriter strike-out}{ega25_25_119-v0011}
\var{\xout{\unsure{considérée comme}}}{}{typewriter strike-out}{ega25_25_119-v0012}  Soit $\emptyset$ l'ensemble des  \var{\deleted{``applications rationelles}}{}{deletion}{ega25_25_119-v0013}\var{\inserted{pseudo-morphismes}}{pseudo-morphismes}{insertion}{ega25_25_119-v0014} rel à
$S$ \var{\deleted{\unsure{*}}}{}{deletion}{ega25_25_119-v0015}\var{\inserted{$f$}}{$f$}{insertion}{ega25_25_119-v0016} de $X$ dans  $P$, telles que le\var{\xout{\illegible{1}}}{}{typewriter strike-out}{ega25_25_119-v0017} domaine de définition \var{\inserted{$U=$}}{$U=$}{insertion}{ega25_25_119-v0018}\var{\over{dom}{\unsure{U}}}{\unsure{U}}{writing over}{ega25_25_119-v0019}(f)
\var{\xout{\unsure{f satisf}}}{}{typewriter strike-out}{ega25_25_119-v0020} satisfasse la condition : $z \in X$-$U$\var{\deleted{(f)}}{}{deletion}{ega25_25_119-v0021} $\Longrightarrow$ prof$_z \underline{O}_{X_s} \geqslant 2$\var{\inserted{$\surd$, et $f_U = f|U : U \to \check{P}$ satisfait la condition de non dégénérescence envisagée dans 3.1.}}{$\surd$, et $f_U = f|U : U \to \check{P}$ satisfait la condition de non dégénérescence envisagée dans 3.1.}{insertion}{ega25_25_119-v0022}. Considé\var{\xout{rons}}{}{typewriter strike-out}{ega25_25_119-v0023}
rons l'application naturelle  $D \rightsquigarrow f_D$  de $\emptyset$  dans $\emptyset$ .  Alors :

a)  Cette application est injective,  et  pour $D \in \emptyset$ , l'ensemble des
points fixes $Z$  n'est autre que l'ensemble  complémentaire du domaine de
définition \var{\xout{\unsure{Dif}}}{}{typewriter strike-out}{ega25_25_119-v0024} $U$\var{\deleted{(f)}}{}{deletion}{ega25_25_119-v0025} de $f$\hand{$_D$}.

b)  \var{\xout{\unsure{Pour qu'une f $\emptyset$ appartienne à provienne d'une $D$ $\emptyset$, il}}}{}{typewriter strike-out}{ega25_25_119-v0026} Soit
$f \in \emptyset$ , et soit $U$ un ouvert sur lequel $f$ soit défini, et tel que $z \in X$-$U$
implique prof $\underline{O}_{X_s,z} \geqslant 2$ \var{\xout{\unsure{et existe}}}{}{typewriter strike-out}{ega25_25_119-v0027} (par exemple $U = U(f)$ domaine de
définition de $f$). Pour que $f$ provienne d'une $D \in \emptyset$ , il f et suffit  que\hand{\unsure{$\surd$}}
\var{\deleted{\unsure{$f_U$ satisfasse la condition énoncée dans 3.1. et que}}}{}{deletion}{ega25_25_119-v0028}
posant $\underline{L}_U = f_U^*(\underline{O}_P(1))$  (où  $f_U$: $U \to P$  est le morphisme induit par $f$),
\var{\deleted{\unsure{$f_U$ satisfasse la condition énoncée dans 3.1 et}}}{}{deletion}{ega25_25_119-v0029}
\var{\xout{\unsure{que}}}{}{typewriter strike-out}{ega25_25_119-v0030} le Module  $i_*(\underline{L}_U)$  sur $X$ soit inversible (où $i$:$U \to X$ est l'immer-%
sion canonique).

On remarquera que lorsque  les fibres de $X$ sur $S$ satisfont à $(S_2)$,
par exemple sont normales, voire géométriquement  normales, alors la
condition de profondeur  envisagée sur  une partie fermée $Z$ de $X$ dans \var{\deleted{\illegible{1}}}{}{deletion}{ega25_25_119-v0031}
la proposition  signifie simplement que pour tout $s \in S$,

\mspage{119}
